% ═══════════════════════════════════════════════════════════════
% MT-02d: Minitest Wiederholung Familie
% Spanisch Klasse 7
% ═══════════════════════════════════════════════════════════════
% Dauer: 15 Minuten
% Punkte: 30
% Inhalt: Familienmitglieder, tener, Possessivbegleiter, Zahlen, Dialog
% ═══════════════════════════════════════════════════════════════

\documentclass[11pt, a4paper]{article}

\usepackage[utf8]{inputenc}
\usepackage[T1]{fontenc}
\usepackage[ngerman]{babel}
\usepackage{helvet}
\renewcommand{\familydefault}{\sfdefault}

\usepackage[margin=2cm]{geometry}
\usepackage{enumitem}
\usepackage{tabularx}
\usepackage{booktabs}

\usepackage{xcolor}
\usepackage{tcolorbox}
\tcbuselibrary{skins}

\usepackage{fancyhdr}
\usepackage{lastpage}
\usepackage{needspace}

% FARBEN
\definecolor{spanisch-color}{HTML}{c41e3a}
\definecolor{grau-color}{HTML}{888888}
\definecolor{hellgrau-color}{HTML}{f5f5f5}

\colorlet{spanisch}{spanisch-color}
\colorlet{grau}{grau-color}
\colorlet{hellgrau}{hellgrau-color}

\newcommand{\fachfarbe}{spanisch}

\newcommand{\thema}{Minitest: Wiederholung Mi familia}
\newcommand{\fach}{Spanisch}
\newcommand{\klasse}{7}
\newcommand{\maxpunkte}{30}
\newcommand{\zeit}{15 Minuten}
\newcommand{\datum}{\today}

\pagestyle{fancy}
\fancyhf{}
\renewcommand{\headrulewidth}{0.4pt}
\renewcommand{\footrulewidth}{0.4pt}

\lhead{\fach{} -- Klasse \klasse}
\chead{\textbf{MT-02d}}
\rhead{\datum}

\lfoot{\zeit}
\cfoot{}
\rfoot{Seite \thepage\ von \pageref{LastPage}}

\newcommand{\punktebox}[1]{\hfill\fbox{\small \_/#1}}
\newcommand{\luecke}[1]{\rule{#1}{0.4pt}}

\newtcolorbox{infobox}{
    colback=hellgrau,
    colframe=\fachfarbe,
    boxrule=1pt,
    arc=0pt,
    left=8pt, right=8pt, top=8pt, bottom=8pt
}

\newenvironment{aufgabe}[1]{%
    \needspace{5\baselineskip}%
    \nopagebreak[4]%
    \vspace{10pt}
    \noindent\textbf{\textcolor{\fachfarbe}{Aufgabe #1}}
    \nopagebreak[4]%
    \vspace{5pt}
    \nopagebreak[4]%
    \begin{list}{}{\leftmargin=0pt}
    \item[]
}{%
    \end{list}
}

\newcommand{\es}[1]{\textcolor{\fachfarbe}{\textbf{#1}}}

\begin{document}

% ═══════════════════════════════════════════════════════════════
% KOPFBEREICH
% ═══════════════════════════════════════════════════════════════

\begin{center}
    {\Large\bfseries\textcolor{\fachfarbe}{\thema}}
\end{center}

\vspace{5pt}

\noindent
\begin{tabularx}{\textwidth}{|X|X|X|}
\hline
\textbf{Name:} \luecke{4cm} &
\textbf{Klasse:} \klasse &
\textbf{Datum:} \luecke{2cm} \\
\hline
\end{tabularx}

\vspace{10pt}

\begin{infobox}
\textbf{Zeit:} \zeit{} | \textbf{Punkte:} \maxpunkte{} | \textbf{Hilfsmittel:} keine
\end{infobox}

\vspace{15pt}

% ═══════════════════════════════════════════════════════════════
% AUFGABE 1: Familienmitglieder zuordnen (6 Punkte)
% ═══════════════════════════════════════════════════════════════

\begin{aufgabe}{1 -- Familienmitglieder zuordnen}
Verbinde die spanischen Familienbegriffe mit der deutschen Bedeutung. \punktebox{6}
\end{aufgabe}

\begin{center}
\begin{tabular}{rcl}
\es{la abuela} & \luecke{2cm} & der Bruder \\[0.8em]
\es{el hermano} & \luecke{2cm} & die Cousine \\[0.8em]
\es{la prima} & \luecke{2cm} & die Großmutter \\[0.8em]
\es{el tío} & \luecke{2cm} & der Onkel \\[0.8em]
\es{la madre} & \luecke{2cm} & der Großvater \\[0.8em]
\es{el abuelo} & \luecke{2cm} & die Mutter \\
\end{tabular}
\end{center}

\vspace{15pt}

% ═══════════════════════════════════════════════════════════════
% AUFGABE 2: tener einsetzen (6 Punkte)
% ═══════════════════════════════════════════════════════════════

\begin{aufgabe}{2 -- tener konjugieren}
Ergänze die richtige Form von \es{tener}. \punktebox{6}
\end{aufgabe}

\noindent
a) Yo \luecke{2.5cm} dos hermanas. \hfill (ich habe)

\vspace{0.8em}

\noindent
b) María \luecke{2.5cm} un perro. \hfill (María hat)

\vspace{0.8em}

\noindent
c) ¿Tú \luecke{2.5cm} primos? \hfill (hast du?)

\vspace{15pt}

% ═══════════════════════════════════════════════════════════════
% AUFGABE 3: Possessivbegleiter (6 Punkte)
% ═══════════════════════════════════════════════════════════════

\begin{aufgabe}{3 -- Possessivbegleiter}
Ergänze: \es{mi}, \es{tu}, \es{su}, \es{mis}, \es{tus} oder \es{sus}. \punktebox{6}
\end{aufgabe}

\noindent
a) \luecke{2cm} madre se llama Carmen. \hfill (meine Mutter)

\vspace{0.8em}

\noindent
b) ¿Cómo se llama \luecke{2cm} padre? \hfill (dein Vater)

\vspace{0.8em}

\noindent
c) \luecke{2cm} hermanos son muy simpáticos. \hfill (seine Brüder)

\vspace{15pt}

% ═══════════════════════════════════════════════════════════════
% AUFGABE 4: Zahlen und Alter (6 Punkte)
% ═══════════════════════════════════════════════════════════════

\begin{aufgabe}{4 -- Alter angeben}
Schreibe die Zahl als spanisches Wort. \punktebox{6}
\end{aufgabe}

\noindent
a) Mi hermano tiene \luecke{3.5cm} años. \hfill (12)

\vspace{0.8em}

\noindent
b) Tengo \luecke{3.5cm} años. \hfill (15)

\vspace{0.8em}

\noindent
c) Mi abuela tiene \luecke{3.5cm} años. \hfill (18)

\vspace{15pt}

% ═══════════════════════════════════════════════════════════════
% AUFGABE 5: Dialog (6 Punkte)
% ═══════════════════════════════════════════════════════════════

\begin{aufgabe}{5 -- Dialog vervollständigen}
Ergänze den Dialog sinnvoll auf Spanisch. \punktebox{6}
\end{aufgabe}

\noindent
\textbf{A:} ¡Hola! ¿Tienes hermanos?

\vspace{0.8em}

\noindent
\textbf{B:} \luecke{8cm}

\vspace{0.8em}

\noindent
\textbf{A:} ¿Cómo se llama tu madre?

\vspace{0.8em}

\noindent
\textbf{B:} \luecke{8cm}

\vspace{0.8em}

\noindent
\textbf{A:} ¿Cuántos años tiene tu padre?

\vspace{0.8em}

\noindent
\textbf{B:} \luecke{8cm}

% ═══════════════════════════════════════════════════════════════
% PUNKTEÜBERSICHT
% ═══════════════════════════════════════════════════════════════

\vspace{25pt}
\hrule
\vspace{10pt}

\noindent
\begin{tabularx}{\textwidth}{|X|c|c|c|c|c||c|}
\hline
\textbf{Aufgabe} & 1 & 2 & 3 & 4 & 5 & \textbf{Gesamt} \\
\hline
\textbf{max. Punkte} & 6 & 6 & 6 & 6 & 6 & \textbf{30} \\
\hline
\textbf{erreicht} & & & & & & \\[0.8em]
\hline
\end{tabularx}

\vspace{10pt}

\begin{flushright}
\textbf{Note:} \luecke{2cm}
\end{flushright}

\end{document}
